\documentclass{article}
\usepackage{mathtools} 
\usepackage{fontspec}
\usepackage[UTF8]{ctex}
\usepackage{amsthm}
\usepackage{mdframed}
\usepackage{xcolor}
\usepackage{amssymb}
\usepackage{amsmath}


% 定义新的带灰色背景的说明环境 zremark
\newmdtheoremenv[
  backgroundcolor=gray!10,
  % 边框与背景一致，边框线会消失
  linecolor=gray!10
]{zremark}{说明}

% 通用矩阵命令: \flexmatrix{矩阵名}{元素符号}{行数}{列数}
\newcommand{\flexmatrix}[4]{
  \[
  #1 = \begin{pmatrix}
    #2_{11}     & #2_{12}     & \cdots & #2_{1#4}   \\
    #2_{21}     & #2_{22}     & \cdots & #2_{2#4}   \\
    \vdots      & \vdots      & \ddots & \vdots     \\
    #2_{#31}    & #2_{#32}    & \cdots & #2_{#3#4}
  \end{pmatrix}
  \]
}

% 简化版命令（默认矩阵名为A，元素符号为a）: \quickmatrix{行数}{列数}
\newcommand{\quickmatrix}[2]{\flexmatrix{A}{a}{#1}{#2}}

\begin{document}
\title{习题 7.2}
\author{张志聪}
\maketitle

\section*{5}

以上极限为例。

由定理7.4可知，有界数列$\{a_n\}, \{b_n\}$都有上极限（即最大聚点），
分别设为$\overline{A}, \overline{B}$。

反证法，假设$\overline{A} > \overline{B}$。
令$\epsilon = \frac{\overline{A} - \overline{B}}{3}$，
于是的
\begin{align*}
  \overline{B} + \epsilon < \overline{A} - \epsilon
\end{align*}
由定理7.7可知，存在$\{a_n\}$的子列$\{a_{n_k}\}$，$\overline{A} + \epsilon > a_{n_k} > \overline{A} - \epsilon$。
同理，存在$\{b_n\}$的子列$\{b_{n_k}\}$，$\overline{B} + \epsilon > b_{n_k} > \overline{B} - \epsilon$。

于是，对任意$N$，都有存在$n > N$有$a_n$是子列$\{a_{n_k}\}$中的项，
同理，存在$n > N$有$b_n$是子列$\{b_{n_k}\}$中的项。
这与存在$N_0 > 0$，当$n > N_0$时，有$a_n \leq b_n$矛盾。

\section*{6}

只证明上极限。

\begin{itemize}
  \item 必要性；

        对任意$\epsilon > 0$，
        由于$\overline{A}$是上极限，由定理7.7(i)可知，
        可知，存在$N_0$，使得$n > N_0$，我们有
        \begin{align*}
          \sup\limits_{k \geq n} \{a_k\} < \overline{A} + \epsilon
        \end{align*}

        由定理7.7(ii)可知，存在$n > n_1$，我们有
        \begin{align*}
          \overline{A} - \epsilon < a_n
        \end{align*}
        所以，
        \begin{align*}
          \sup\limits_{k \geq n_1} \{a_k\} \geq \overline{A} - \epsilon
        \end{align*}

        综上，取$N = max(N_0, n_1)$，对$n > N$，我们有
        \begin{align*}
          \overline{A} - \epsilon \leq \sup\limits_{k \geq n} \{a_k\} < \overline{A} + \epsilon
        \end{align*}
        所以
        \begin{align*}
          \lim \limits_{n \to \infty} \sup\limits_{k \geq n} \{a_k\} = \overline{A}
        \end{align*}

  \item 充分性

        (i)

        由极限定义可知，对任意$\epsilon>0$，存在$N > 0$，使得对任意$n > N$，有
        \begin{align*}
          |\sup\limits_{k \geq n} \{a_k\} - \overline{A}| < \epsilon \\
          \overline{A} - \epsilon < \sup\limits_{k \geq n} \{a_k\} < \overline{A} + \epsilon
        \end{align*}
        由上确界的定义可知，对$n > N$，都有
        \begin{align*}
          a_n < \sup\limits_{k \geq n} \{a_k\}
        \end{align*}
        即
        \begin{align*}
          a_n < \overline{A} + \epsilon
        \end{align*}

        在再次利用上确界定义可知，存在$n_0 > N$，有
        \begin{align*}
          \overline{A} - \epsilon < a_{n_0} < \sup\limits_{k \geq n} \{a_k\}
        \end{align*}
        又因为，对任意$n > n_0 (> N)$，我们也有
        \begin{align*}
          |\sup\limits_{k \geq n} \{a_k\} - \overline{A}| < \epsilon
        \end{align*}
        类似的操作，存在$n_1 > n_0$，有
        \begin{align*}
          \overline{A} - \epsilon < a_{n_1} < \sup\limits_{k \geq n} \{a_k\}
        \end{align*}
        重复操作，可得子列$\{a_{n_k}\}, a_{n_k} > \overline{A} - \epsilon$。

        综上，利用定理7.7可知，$\overline{A}$是上极限。

\end{itemize}


\end{document}